%% STEMKit -- SoftwareX (Original Software Publication) manuscript
%%
%% Prepared with the SoftwareX article template (elsarticle, numbered
%% references). Structure follows SoftwareX article template Version 6
%% (March 2026): frontmatter, code metadata table, and the five mandatory
%% sections 1-5, followed by Acknowledgements and References.
%%
%% Build:  pdflatex preprint   (run twice, for cross-references)
%%
%% Copyright 2026 Elsevier Ltd for the template; distributed under the
%% conditions of the LaTeX Project Public License, version 1.2 or later,
%% http://www.latex-project.org/lppl.txt

\documentclass[preprint,12pt,a4paper]{elsarticle}

%% Use the option review to obtain double line spacing
%% \documentclass[authoryear,preprint,review,12pt]{elsarticle}

%% For including figures, graphicx.sty has been loaded in elsarticle.cls.

\usepackage{amsmath}
\usepackage{amssymb}
\usepackage{booktabs}
\usepackage{siunitx}
\usepackage{xcolor}
\usepackage{listings}
\usepackage{hyperref}
%% Long \texttt{} strings such as file paths and package names cannot hyphenate
%% and would otherwise overflow the margin; this permits breaks inside them.
\usepackage[htt]{hyphenat}

%% The template sets flush-left paragraphs; a small vertical skip is added so
%% that consecutive paragraphs remain distinguishable.
\setlength{\parindent}{0pt}
\setlength{\parskip}{6pt}

%% The lineno package adds line numbers. Start line numbering with
%% \begin{linenumbers}, end it with \end{linenumbers}. Or switch it on for the
%% whole article with \linenumbers.
%\usepackage{lineno}

\sisetup{detect-all}

%% Listing style for the code excerpts.
\definecolor{codebg}{RGB}{248,249,250}
\definecolor{codekw}{RGB}{0,92,197}
\definecolor{codecm}{RGB}{106,115,125}
\definecolor{codestr}{RGB}{3,47,98}

\lstdefinelanguage{JavaScriptES6}{
  keywords={import, from, export, const, let, function, return, await, async,
            if, else, for, of, new, class, extends, throw, try, catch},
  keywordstyle=\color{codekw}\bfseries,
  ndkeywords={true, false, null, undefined, NaN, Infinity},
  ndkeywordstyle=\color{codekw},
  comment=[l]{//},
  morecomment=[s]{/*}{*/},
  commentstyle=\color{codecm}\itshape,
  stringstyle=\color{codestr},
  morestring=[b]',
  morestring=[b]",
  morestring=[b]`,
  sensitive=true
}

\lstset{
  language=JavaScriptES6,
  backgroundcolor=\color{codebg},
  basicstyle=\ttfamily\footnotesize,
  breaklines=true,
  captionpos=b,
  frame=single,
  framesep=4pt,
  rulecolor=\color{gray!30},
  showstringspaces=false,
  xleftmargin=6pt,
  xrightmargin=6pt
}

%% elsarticle's preprint style prints "Preprint submitted to <journal>" and the
%% date in the footer of every page. This restores a plain centred page number.
\makeatletter
\def\ps@pprintTitle{%
  \let\@oddhead\@empty
  \let\@evenhead\@empty
  \def\@oddfoot{\reset@font\hfil\thepage\hfil}%
  \let\@evenfoot\@oddfoot}
\makeatother

\begin{document}
\renewcommand{\labelenumii}{\arabic{enumi}.\arabic{enumii}}

\begin{frontmatter}

%% Title, authors and addresses

\title{STEMKit: A client-side toolkit for computational chemistry and
scientific data analysis}

\author[1]{Olanrewaju M. Daramola\corref{cor1}\fnref{orcid}}
\ead{lanrelangmuir@gmail.com}
\cortext[cor1]{Corresponding author.}
\fntext[orcid]{ORCID:
\href{https://orcid.org/0009-0006-3327-2047}{0009-0006-3327-2047}}

\address[1]{Independent Researcher}

\begin{abstract}
%% Text of abstract: ca. 100 words
STEMKit is a suite of 21 browser-based tools for computational chemistry and
scientific data analysis, built on \texttt{@stemkit/core}, a JavaScript library
of parsing and numerical routines carrying no dependency on the document object
model. All computation executes inside the user's own browser: no data is
transmitted, no account is required, and nothing is installed, so the tools
remain usable for unpublished or confidential material. The identical modules
also execute under Node.js, so an analysis explored interactively can be
captured as a version-pinned script and placed under continuous integration.
Sixteen domain modules are covered by 1075 tests validated against SciPy,
NumPy,
and physical invariants rather than against the implementation itself.
\end{abstract}

\begin{keyword}
%% keywords here, in the form: keyword \sep keyword
Computational chemistry \sep Molecular dynamics \sep Client-side computation
\sep GROMACS \sep Reproducibility \sep Data privacy
\end{keyword}

\end{frontmatter}

%\linenumbers

% ---------------------------------------------------------------------------
\section*{Required Metadata}

\section*{Current code version}

\begin{table}[!h]
\begin{tabular}{|l|p{4.4cm}|p{6.3cm}|}
\hline
\textbf{Nr.} & \textbf{Code metadata description} & \textbf{Metadata} \\
\hline
C1 & Current code version & v0.1.0 \\
\hline
C2 & Permanent link to code/repository used for this code version &
\url{https://doi.org/10.5281/zenodo.21543112} (archived release \texttt{v0.1.0});
repository at \url{https://github.com/LD-Shell/stemkit} \\
\hline
C3 & Legal code license & MIT License \\
\hline
C4 & Code versioning system used & git \\
\hline
C5 & Software code languages, tools and services used & JavaScript
(ECMAScript 2020 modules), HTML5, CSS3; Node.js; Jest; Tailwind CSS; GitHub
Actions \\
\hline
C6 & Compilation requirements, operating environments and dependencies &
Browser tools: any current browser supporting ECMAScript modules and the
\texttt{FileReader} API; no installation and no build step. Library: Node.js
$\geq 18$; \texttt{npm install} installs the development dependencies (Jest,
Tailwind CSS) only. There are no runtime dependencies --- jStat, Papa Parse,
regression.js and bibtex-parse-js are vendored in the repository \\
\hline
C7 & If available, link to developer documentation/manual &
\url{https://github.com/LD-Shell/stemkit#readme}; per-module API documentation
in \texttt{src/core/README.md}; hosted tools at \url{https://stemkit.net} \\
\hline
C8 & Support email for questions & lanrelangmuir@gmail.com (issue tracker:
\url{https://github.com/LD-Shell/stemkit/issues}) \\
\hline
\end{tabular}
\caption{Code metadata (mandatory)}
\label{tab:metadata}
\end{table}

% ---------------------------------------------------------------------------
\section{Motivation and significance}
\label{sec:motivation}

The software landscape of computational chemistry is well served at its
extremes. Simulation engines such as GROMACS~\cite{abraham2015} and
LAMMPS~\cite{thompson2022} are mature, extensively validated, and actively
maintained; analysis frameworks such as MDAnalysis~\cite{michaudagrawal2011}
provide comprehensive programmatic access to trajectory data. Between these
poles lies a substantial volume of routine work that neither category serves
comfortably.

Consider the interval between completing a simulation and reporting it. A
researcher must inspect the RMSD to confirm equilibration, extract a radius of
gyration, convert a structure between the nanometre convention of GROMACS and
the \si{\angstrom} convention of the Protein Data Bank, compare replicate
measurements statistically, digitise a figure from a paper for comparison,
abbreviate journal titles in a reference list, and compose a scheduler script
requesting an appropriate resource allocation. None of these operations is
intellectually demanding. All of them are error-prone when performed by hand,
and all of them are performed repeatedly.

\subsection{The friction problem}

Three classes of solution exist, and each imposes a characteristic cost.

\emph{Local scripting.} A Python environment with MDAnalysis, NumPy and SciPy
addresses every task listed above and is the correct instrument for sustained
analytical work. Its cost is environmental. Managed institutional workstations
frequently prohibit package installation; conda environments drift; and a
script written six months ago may no longer execute against the current
interpreter. For a researcher whose primary expertise lies in chemistry rather
than software engineering, this cost is not negligible.

\emph{Web services.} Server-side tools eliminate installation entirely but
require data transmission. For unpublished results, material subject to a
confidentiality agreement, or work under embargo, this is frequently prohibited
outright. Even where permitted, the transfer creates a record whose retention
policy the researcher does not control.

\emph{Desktop applications.} Graphical applications avoid both problems but
introduce a third: they are difficult to script. An analysis performed by
clicking cannot readily be re-executed, version-controlled, or subjected to
continuous integration. In practice such an analysis is difficult to reproduce,
even when every step was recorded at the time.

\subsection{The client-side position}

Modern browsers provide a capable computational environment. The
\texttt{FileReader} API grants access to local files without transmission, and
JavaScript engines execute numerical code fast enough for the data volumes
involved. ECMAScript modules make it possible to organise that code properly,
instead of accumulating scripts.

STEMKit is built around these constraints. No server participates in the
computation, so there is nothing to which data could be transmitted and users
are not asked to take a privacy policy on trust. The reproducibility property
derives from a separate architectural decision: the numerical core is extracted
into modules with no dependency on the browser, so the identical code that
executes behind the interface also executes under Node.js.

\subsection{Experimental setting}

The software is used in two ways, which share one implementation.
Interactively, the user opens a tool page at \url{https://stemkit.net} (or a
local clone) and selects a local file --- a GROMACS \texttt{.xvg} series, a PDB
or \texttt{.gro} structure, a CSV of replicate measurements --- and reads the
result in the page. The file is read through \texttt{FileReader}; it is never
uploaded. Programmatically, the same routines are imported from
\texttt{@stemkit/core} in a Node.js script, which is the path taken once an
exploratory analysis is to be fixed, version-pinned and re-run. Every tool that
produces a figure additionally emits a standalone matplotlib script that
regenerates that figure from the original data file, so an interactive result
can be reproduced offline without the browser.

Related work is cited where the corresponding functionality is described in
Section~\ref{sec:description}: GROMACS~\cite{abraham2015} and
LAMMPS~\cite{thompson2022} for the trajectory and job-script formats,
PLUMED~\cite{tribello2014} for collective-variable output,
MDAnalysis~\cite{michaudagrawal2011} as the reference framework for trajectory
analysis, SciPy~\cite{virtanen2020} and NumPy~\cite{harris2020} as the
validation references for the numerics, and the primary statistical
literature~\cite{welch1947,dagostino1971,brown1974,holm1979,iglewicz1993,grubbs1969}
for the implemented tests. Conversion factors follow CODATA
2018~\cite{tiesinga2021}, and the numerical treatment of tail probabilities
follows standard practice for the incomplete beta and gamma
functions~\cite{press2007}.

% ---------------------------------------------------------------------------
\section{Software description}
\label{sec:description}

\subsection{Software architecture}
\label{sec:architecture}

The original implementation followed the pattern common to browser tooling:
each of the 21 tools consisted of an HTML document paired with a single
JavaScript file combining event handling, rendering and computation. This
arrangement is expedient during development and progressively obstructive
thereafter. Computational logic entangled with DOM manipulation cannot be
tested without a browser, cannot be reused across tools, and cannot be invoked
from a script.

The refactoring applied the \emph{strangler fig} pattern: for each tool, the
mathematical and parsing logic was extracted into a pure module beneath
\texttt{src/core/}, a test suite was written against the extracted module, and
the original file was then reduced to DOM wiring that imports the new
functions. The user-visible interface is unchanged; the computation beneath it
is now independently addressable. Figure~\ref{fig:architecture} shows the
resulting layering and Table~\ref{tab:modules} enumerates the modules.

\begin{figure}[htbp]
\centering
\small
\begin{tabular}{c}
\begin{tabular}{@{}c@{\hspace{1.2em}}c@{}}
\fbox{\parbox[t]{0.43\textwidth}{\centering \textbf{Browser host}\\[2pt]
21 static HTML/CSS pages\\ DOM wiring only; \texttt{FileReader} input\\
UMD bundles loaded by \texttt{<script>}}} &
\fbox{\parbox[t]{0.43\textwidth}{\centering \textbf{Node.js host}\\[2pt]
user analysis scripts\\ 1075-test Jest suite; smoke test\\
UMD bundles via \texttt{createRequire}}} \\
\end{tabular} \\[6pt]
$\downarrow$ \hspace{0.4\textwidth} $\downarrow$ \\[6pt]
\fbox{\parbox{0.9\textwidth}{\centering \textbf{\texttt{@stemkit/core}}\\[2pt]
16 DOM-free domain modules, aggregated by \texttt{src/core/index.js}}} \\[6pt]
$\downarrow$ \\[6pt]
\fbox{\parbox{0.9\textwidth}{\centering \textbf{Vendor injection layer}\\[2pt]
\texttt{registerVendor()} (Node.js) $\;\cdot\;$
\texttt{registerFromGlobals()} (browser)}} \\[6pt]
$\downarrow$ \\[6pt]
\fbox{\parbox{0.9\textwidth}{\centering \textbf{Vendored UMD bundles}\\[2pt]
jStat $\;\cdot\;$ Papa Parse $\;\cdot\;$ regression.js $\;\cdot\;$
bibtex-parse-js}} \\
\end{tabular}
\caption{Architecture of STEMKit. Both hosts execute the identical core
modules; the only environment-specific code is the registration call that
supplies the vendored libraries.}
\label{fig:architecture}
\end{figure}

\begin{table}[htbp]
\centering
\caption{Modules of \texttt{@stemkit/core}. Vendored dependencies are supplied
by injection; modules marked \emph{none} operate without any third-party code.}
\label{tab:modules}
\footnotesize
\begin{tabular}{@{}llrl@{}}
\toprule
\textbf{Module} & \textbf{Domain} & \textbf{Tests} & \textbf{Dependency} \\
\midrule
\texttt{bibtex} & Reference deduplication, sanitising & 102 & bibtex-parse-js \\
\texttt{structure} & Molecular geometry, PDB/GRO/XYZ & 100 & none \\
\texttt{statistics} & Inferential statistics & 82 & jStat \\
\texttt{curve-fitting} & Least-squares regression & 77 & regression.js \\
\texttt{xvg-parser} & GROMACS/PLUMED trajectory data & 77 & none \\
\texttt{units} & Physical unit conversion & 73 & none \\
\texttt{plumed} & PLUMED input generation & 72 & none \\
\texttt{selection} & Atom selection, spatial queries & 67 & none \\
\texttt{slurm} & HPC job-script generation & 61 & none \\
\texttt{latex} & Table generation and escaping & 59 & none \\
\texttt{error-bars} & Group summaries, error bars & 57 & jStat \\
\texttt{journals} & Whole-title journal abbreviation & 53 & none \\
\texttt{data-cleaning} & Tabular transformation & 51 & Papa Parse \\
\texttt{outliers} & Anomaly detection & 51 & jStat \\
\texttt{digitizer} & Figure digitisation & 49 & none \\
\texttt{iso4} & ISO 4 word-level abbreviation (LTWA) & 44 & none \\
\midrule
\textbf{Total} & & \textbf{1075} & \\
\bottomrule
\end{tabular}
\end{table}

Four third-party libraries are vendored within the repository: jStat for
statistical distributions, Papa Parse for delimited-text parsing, regression.js
for curve fitting, and bibtex-parse-js for reference parsing. Each is
distributed as a Universal Module Definition (UMD) bundle, which detects its
host at run time and either assigns \texttt{module.exports} under CommonJS or
attaches a global to \texttt{window} in a browser. Neither path is reachable
from a plain ECMAScript module: a UMD file contains no \texttt{export}
statements, so a direct import yields no bindings, while
\texttt{import \{ createRequire \} from 'module'} --- the conventional Node.js
workaround --- constitutes a hard \emph{resolution} failure in a browser, which
occurs before any code executes and therefore cannot be intercepted by
\texttt{try}/\texttt{catch}. Rather than branch on the execution environment
inside every module, which would render the core untestable in one of its two
targets, the library declares the libraries it requires and each host registers
them at initialisation (Section~\ref{sec:snippets}). Modules request their
dependencies lazily, at the point of use rather than at import time, so
importing any module never fails merely because an unrelated library has not yet
been registered.

One consequence merits explicit note, as it is a common failure mode. A
repository declaring \texttt{"type": "module"} instructs Node.js to parse every
\texttt{.js} file beneath it as an ECMAScript module, including the vendored UMD
bundles; under that interpretation the UMD factory takes its browser branch and
fails. Scoping \texttt{"type": "commonjs"} to the dependency directory alone
restores correct behaviour.

The test suite comprises 1075 tests across the 16 domain modules, with 93.9\%
statement and 97.6\% line coverage of \texttt{src/core/}. Its governing
principle is that numerical results are validated against \emph{independent}
references rather than against the implementation under test, since a test
written from the same source as the code confirms only internal consistency.
Reference values were obtained from SciPy~\cite{virtanen2020} for t-tests,
ANOVA, correlation, non-parametric tests, quantiles and multiple-comparison
correction; from NumPy~\cite{harris2020} for regression coefficients and
descriptive statistics; from \texttt{scipy.constants} for every unit conversion
factor; and from physical invariants for structural geometry --- specifically
the molecular weight and centre of mass of water, the orthonormality of
rotation matrices, the preservation of interatomic distances under rotation, and
round-trip fidelity through every supported file format.
An end-to-end smoke
test additionally exercises one representative path through fifteen of the
sixteen modules against a real installation, detecting the class of failure that
unit tests cannot: a broken aggregate export, a misconfigured module type
declaration, or a vendored bundle that fails to load.

\subsection{Software functionalities}
\label{sec:functionalities}

\subsubsection{Trajectory and collective-variable data}

GROMACS writes analysis output in the Grace format, a plain-text convention in
which lines beginning with \texttt{\#} are comments, lines beginning with
\texttt{@} are formatting directives, and all remaining non-empty lines are
whitespace-delimited numeric records. The directives carry the plot title, axis
labels and per-series legends. Discarding them --- as a naive parser that skips
all non-numeric lines would --- loses precisely the metadata that identifies
which column contains which observable. The \texttt{xvg-parser} module recovers
this metadata alongside the numeric matrix. Malformed records are counted and
reported, not silently discarded, so a truncated trajectory is visible to the
user instead of quietly shortening the analysis. Two further details of the
format are handled explicitly: the Grace multi-set separator \texttt{\&} is
treated as a delimiter rather than parsed as data, and records containing
\texttt{NaN} or \texttt{Infinity} are rejected whole, since one non-finite entry
would otherwise contaminate every downstream statistic.

PLUMED~\cite{tribello2014} writes \texttt{COLVAR} files whose header lines begin
with \texttt{\#!} and carry field names and metadata. Although the convention
differs from Grace, the comment character coincides, and the same parser
therefore handles both formats without modification. This is verified by an
explicit test rather than left to assumption.

\subsubsection{Structure files and molecular geometry}

The \texttt{structure} module parses the three coordinate formats in routine use
for classical simulation: PDB and GROMACS \texttt{.gro}, both fixed-column
formats in which fields are located by character position, and XYZ, which is
whitespace-delimited. Fixed-column parsing is essential rather than pedantic: a
residue name may legitimately be blank, and a whitespace-splitting parser would
silently shift every subsequent field on such a line. Unit handling is explicit
throughout --- PDB and XYZ express coordinates in \si{\angstrom}, \texttt{.gro}
uses nanometres, and every parse result carries its unit --- because silently
mixing the two is the most direct route to a structure that is wrong by a factor
of ten.

Assigning an element to each atom is a prerequisite for any mass-weighted
quantity, and PDB atom names are ambiguous by construction: \texttt{CA} denotes
the $\alpha$-carbon in a protein residue and calcium in an ion record. The
resolution implemented here proceeds in five stages.

\begin{enumerate}
  \item An explicit element column, where present, takes priority.
  \item A leading digit is a hydrogen-count prefix (\texttt{1HB},
        \texttt{2HG1}) and is stripped before the symbol is read.
  \item An unambiguous two-letter metal (\texttt{FE}, \texttt{ZN},
        \texttt{SE}) is recognised wherever it appears, including within a host
        residue such as haem or selenomethionine.
  \item Otherwise, a two-letter prefix whose second character is uppercase is a
        PDB remoteness indicator, and only the first letter is the element:
        \texttt{CA} in a protein residue is C$\alpha$.
  \item As a final fallback, a two-letter symbol is trusted when the residue
        shares its name, the convention for monatomic ions.
\end{enumerate}

Stages 2 and 3 correct defects present in the original implementation
(Section~\ref{sec:defects}). The resolution is validated against 40 atom names
drawn from real structures, covering protein backbone and side-chain atoms,
nucleic acid atoms, water, monatomic ions and metal centres. On this basis the
module computes the geometric centroid, the mass-weighted centre of mass, the
axis-aligned bounding box, and the radius of gyration
\begin{equation}
  R_g = \sqrt{\frac{\sum_i m_i \left| \mathbf{r}_i - \mathbf{R} \right|^2}
                   {\sum_i m_i}},
  \qquad
  \mathbf{R} = \frac{\sum_i m_i \mathbf{r}_i}{\sum_i m_i},
\end{equation}
together with rigid-body rotation about an arbitrary pivot using intrinsic
$Z$--$Y$--$X$ Euler angles. Velocities present in a \texttt{.gro} file rotate
with the frame but are never translated, so a round trip through rotation and
export preserves the kinetic state.

\subsubsection{Statistics, outliers and curve fitting}

The \texttt{statistics} module implements Welch's~\cite{welch1947} and Student's
t-tests, the paired t-test, one-way ANOVA, Pearson correlation, the
Mann--Whitney $U$ test and the Wilcoxon signed-rank test. Welch's test is the
default for independent samples, as it does not assume equal variances and loses
negligible power when they are in fact equal. Every parametric test reports an
effect size and, where standard, a confidence interval. Assumptions are checked
and surfaced rather than presumed: normality by the D'Agostino--Pearson $K^2$
omnibus test~\cite{dagostino1971}, and homogeneity of variance by the
median-centred Brown--Forsythe variant of Levene's test~\cite{brown1974}, chosen
for its markedly greater robustness to non-normality than the original
mean-centred formulation. Where an assumption fails, a non-parametric
alternative is recommended. Table~\ref{tab:methods} maps each computed quantity
to its numerical method and validation reference.

\begin{table}[htbp]
\centering
\caption{Numerical methods and validation references.}
\label{tab:methods}
\small
\begin{tabular}{@{}p{3.4cm}p{4.7cm}p{4.5cm}@{}}
\toprule
\textbf{Quantity} & \textbf{Method} & \textbf{Validated against} \\
\midrule
Student / Welch $t$, $p$ &
  Regularised incomplete beta $I_x(a,b)$ &
  \texttt{scipy.stats.ttest\_ind} \\
Welch--Satterthwaite $\nu$ &
  Closed form &
  SciPy, hand calculation \\
One-way ANOVA $F$, $p$ &
  Sum-of-squares decomposition; $I_x$ tail &
  \texttt{scipy.stats.f\_oneway} \\
Pearson $r$, CI &
  Fisher $z$ transform &
  \texttt{scipy.stats.pearsonr} \\
Mann--Whitney $U$ &
  Tie-corrected normal approximation &
  \texttt{scipy.stats.\allowbreak{}mannwhitneyu} \\
Wilcoxon $W$ &
  Signed-rank, zeros discarded &
  \texttt{scipy.stats.wilcoxon} \\
D'Agostino--Pearson $K^2$ &
  Skewness and kurtosis transforms &
  \texttt{scipy.stats.normaltest} \\
Levene $W$ &
  Brown--Forsythe (median-centred) &
  \texttt{scipy.stats.levene} \\
Quantiles &
  Linear interpolation (type 7) &
  \texttt{numpy.percentile} \\
Multiple comparisons &
  Holm--Bonferroni~\cite{holm1979} &
  Independent implementation \\
Grubbs $G$ &
  Inverted $t$ critical value &
  \texttt{scipy.stats.t.ppf} \\
Modified $Z$-score &
  Median/MAD, $0.6745$ scaling~\cite{iglewicz1993} &
  \texttt{numpy} \\
Least-squares fitting &
  Normal equations (regression.js) &
  \texttt{numpy.polyfit} \\
\bottomrule
\end{tabular}
\end{table}

Upper-tail probabilities are computed through the complementary form of the
incomplete beta and gamma functions rather than as $1 - F(x)$. The two
expressions are algebraically identical; numerically they are not. Once $F(x)$
rounds to unity in double precision, the subtraction cancels completely and the
result is exactly zero. The consequence is not academic. For the one-way ANOVA
used in the test suite ($F = 377.545$ on $2$ and $21$ degrees of freedom), the
naive form returns $p = 0$, whereas
\begin{equation}
  p = I_{\frac{d_2}{d_2 + d_1 F}}\!\left(\tfrac{d_2}{2}, \tfrac{d_1}{2}\right)
    = \num{3.4611747382914e-17},
\end{equation}
agreeing with SciPy to twelve significant figures. A continuous test statistic
cannot produce $p = 0$; when that value appears in published output it indicates
a numerical fault in the software, not an unusually strong effect. The same
treatment is applied to the Student $t$, $\chi^2$ and normal tails. Beyond
$|z| \approx 8$ the vendored \texttt{erfc} implementation underflows, and a
documented asymptotic expansion takes over so that a $p$-value is reported as
\num{1.5e-23} rather than as zero.

The \texttt{outliers} module provides Tukey fences, the Iglewicz--Hoaglin
modified $Z$-score~\cite{iglewicz1993} and Grubbs' test~\cite{grubbs1969};
\texttt{curve-fitting} performs least-squares fitting with model-adequacy
warnings. One inherited limitation is documented rather than fixed, and it
applies to two of the three transformed models rather than to all of them. The
vendored regression.js library fits the \emph{power} model by regressing
$\ln y$ on $\ln x$ unweighted, so the residuals minimised are those in
logarithmic space and small $y$ values carry more influence than a direct fit
would give them. The \emph{exponential} model also works in log space but
weights each point by $y$, which is the standard correction for log-transform
bias and keeps the result closer to a direct non-linear fit. The
\emph{logarithmic} model carries none of this bias: it regresses $y$ on
$\ln x$, so only the predictor is transformed and the residuals minimised are
those in $y$, exactly as for any model linear in its parameters. Where the bias
does apply it is material --- for a clean doubling series the weighted
exponential fit returns a growth rate of $\num{0.69022}$ where unweighted
log-space least squares gives $\ln 2 = \num{0.69315}$. Neither value is
incorrect, since they answer different questions, but the difference is large
enough to matter when a rate constant is reported. \texttt{fitCurve} therefore
exposes a \texttt{linearised} flag, true for the exponential and power models
only, so callers may surface the distinction; users requiring publication-grade
non-linear fits are directed to Levenberg--Marquardt optimisation on
untransformed data. A dataset containing points a model cannot transform is
refused with a message naming the requirement --- exponential needs every
$y > 0$, power every $x > 0$ and $y > 0$, logarithmic only $x > 0$ --- rather
than having the offending points silently dropped.

\subsubsection{Digitisation, unit conversion and authoring utilities}

Recovering numerical values from a published figure requires a two-point affine
calibration per axis. On a logarithmic axis the interpolation must be performed
in logarithmic space,
\begin{equation}
  v(p) = 10^{\,\log_{10} v_1
              + \dfrac{(p - p_1)\left(\log_{10} v_2 - \log_{10} v_1\right)}
                      {p_2 - p_1}},
\end{equation}
because a decade occupies a constant pixel distance there. Treating a
logarithmic axis linearly is the most frequent source of error in figure
digitisation, and it produces values that are plausible rather than obviously
wrong. The module additionally reports the data-space uncertainty corresponding
to a one-pixel positioning error, since the precision of a digitised point is
set by the resolution of the figure and not by the number of decimal places the
export happens to print.

Sixty-four units across ten categories are supported, each defined by its ratio
to a category base unit, so conversion reduces to
$x_{\mathrm{to}} = x_{\mathrm{from}} \cdot
f_{\mathrm{to}} / f_{\mathrm{from}}$. This keeps the table linear rather than
quadratic in the number of units and renders every entry independently
checkable against its cited source; factors are CODATA 2018
values~\cite{tiesinga2021}, each verified against \texttt{scipy.constants}
during testing. Temperature is handled separately, since the Celsius and
Fahrenheit scales are affine rather than multiplicative.

Composing a scheduler script correctly requires knowledge that is
engine-specific and easily misapplied, and two aspects dominate the resulting
waste of allocation. First, the resource model differs by engine: GROMACS is
threaded and requires one MPI rank per node with many CPUs per task, whereas
LAMMPS~\cite{thompson2022} is MPI-parallel and requires many ranks per node.
Emitting the wrong shape either idles most of a node or oversubscribes it.
Second, and less widely appreciated, the SLURM \texttt{-{}-mem} directive in an
array job specifies memory \emph{per task}, not per job: a request of 32~GB
across a hundred-task array asks the scheduler for 3.2~TB concurrently. The
generator computes the true in-flight total and reports it. Generated scripts
include checkpoint-resume logic, so a requeued job continues from its last
checkpoint rather than silently restarting from $t = 0$ --- a failure mode that
consumes an entire allocation while appearing to succeed. Warnings are returned
as structured objects rather than formatted strings, leaving presentation to the
caller. The remaining modules support manuscript preparation: LaTeX and Markdown
table generation with correct escaping, ISO-4 journal title abbreviation, and
BibTeX deduplication by union-find over normalised DOIs and titles.

\subsubsection{Defects exposed by extraction}
\label{sec:defects}

Separating computation from presentation subjected the numerical code to
independent verification for the first time, and three defects were identified
that had propagated into reported output. First, skewness and kurtosis are
defined against the \emph{population} standard deviation,
\begin{equation}
  G_1 = \frac{\sqrt{n(n-1)}}{n-2}
        \cdot \frac{1}{n}\sum_{i=1}^{n}
        \left(\frac{x_i - \bar{x}}{\sigma_{\mathrm{pop}}}\right)^{\!3},
  \qquad
  \sigma_{\mathrm{pop}} = \sqrt{\frac{1}{n}\sum_i (x_i - \bar{x})^2},
\end{equation}
whereas the original implementation used the Bessel-corrected sample standard
deviation, deflating skewness by a factor of $\left((n-1)/n\right)^{3/2}$ ---
approximately 15\% at $n = 10$. Because both moments feed the
D'Agostino--Pearson statistic, every normality $p$-value the tool reported was
affected. Second, all upper-tail probabilities were computed by subtraction and
floored at zero for strong effects. Third, haem iron was assigned the mass of
fluorine (18.998~Da rather than 55.845~Da), selenomethionine selenium the mass
of sulfur, and numeric-prefixed hydrogens no mass at all, so molecular weights
and centres of mass computed for metalloproteins were correspondingly wrong.
Each defect is corrected and covered by a regression test. Users who generated
figures with earlier versions of these tools should re-check any reported
skewness, normality $p$-value, or metalloprotein mass.

\subsection{Sample code snippets analysis}
\label{sec:snippets}

Listing~\ref{lst:register} contains the whole of the environment-specific code
in the library. Each host registers the vendored bundles once, after which the
identical modules run unchanged in both.

\begin{lstlisting}[caption={Registration of the vendored libraries under
Node.js and in the browser.},label={lst:register}]
// Node.js: the vendored bundles are loaded through createRequire.
import { createRequire } from 'module';
import { registerVendor } from '@stemkit/core';

const require = createRequire(import.meta.url);
registerVendor({ jStat: require('./js/dependencies/jstat.min.js') });

// Browser: the UMD <script> tags have already installed their globals.
import { registerFromGlobals } from './src/core/index.js';
registerFromGlobals();
\end{lstlisting}

% ---------------------------------------------------------------------------
\section{Illustrative examples}
\label{sec:examples}

The example below is the post-simulation triage that motivated the toolkit:
confirm that a production run has equilibrated, characterise the final
structure, test a replicate comparison and queue a continuation. Every step is
available interactively as a tool page; the scripted form is shown because it is
the reproducible one.

Listing~\ref{lst:xvg} reads a GROMACS RMSD series. The parser returns the
numeric matrix together with the title and per-series legends recovered from the
\texttt{@} directives, so each column is reported under its own observable name
rather than by index.

\begin{lstlisting}[caption={Parsing a GROMACS RMSD trajectory and summarising
each series.},label={lst:xvg}]
import { parseXvg, extractColumn, columnStats } from '@stemkit/core';
import { readFileSync } from 'fs';

const { matrix, headers, title } = parseXvg(readFileSync('rmsd.xvg', 'utf8'));
// title   -> "RMSD & Radius of Gyration"
// headers -> ["Time (ps)", "Backbone RMSD", "Rg"]

headers.slice(1).forEach((name, i) => {
  const s = columnStats(extractColumn(matrix, i + 1));
  console.log(`${name}: ${s.mean.toFixed(4)} +/- ${s.std.toFixed(4)} nm`);
});
\end{lstlisting}

Listing~\ref{lst:structure} takes the final coordinates and computes the
mass-weighted centre of mass and radius of gyration. Because the input is a
\texttt{.gro} file its coordinates are in nanometres, and the conversion to
\si{\angstrom} is requested explicitly rather than assumed.

\begin{lstlisting}[caption={Geometry of the final frame, with an explicit unit
conversion.},label={lst:structure}]
import { parseStructure, centreOfMass, radiusOfGyration, convert }
  from '@stemkit/core';

const frame = parseStructure(readFileSync('final.gro', 'utf8'), 'gro');
const com = centreOfMass(frame.atoms);          // nm, mass-weighted
const rg  = radiusOfGyration(frame.atoms, com); // nm

console.log(`Rg = ${convert(rg, 'nm', 'angstrom').toFixed(2)} A`);
\end{lstlisting}

Listing~\ref{lst:stats} compares two sets of replicate measurements. The
returned object carries the assumption checks alongside the test result, so a
violated assumption is visible at the point of use rather than discovered later;
here a failed normality check redirects the comparison to Mann--Whitney.

\begin{lstlisting}[caption={A replicate comparison with the assumption checks
surfaced.},label={lst:stats}]
import { welchTest, mannWhitneyU } from '@stemkit/core';

const r = welchTest(wildType, mutant);
// r.p, r.effectSize (Hedges' g), r.ci, r.assumptions

if (!r.assumptions.normality.passed) {
  console.warn(r.assumptions.normality.recommendation);
  console.log(mannWhitneyU(wildType, mutant));
} else {
  console.log(`Welch t(${r.df.toFixed(1)}) = ${r.t.toFixed(3)}, p = ${r.p}`);
}
\end{lstlisting}

Finally, Listing~\ref{lst:slurm} generates the submission script for the
continuation run. The engine argument selects the GROMACS resource shape --- one
rank per node, many CPUs per task --- and the returned warnings report the true
in-flight memory total for the array as well as any request the scheduler is
likely to reject.

\begin{lstlisting}[caption={Generating a GROMACS submission script with
resource warnings.},label={lst:slurm}]
import { generateScript } from '@stemkit/core';

const { script, warnings } = generateScript({
  engine: 'gromacs', jobName: 'prod_md', partition: 'gpu',
  nodes: 1, gpus: 1, cpusPerTask: 16,
  walltime: '24:00:00', memory: '32G',
  modules: ['gcc/11.3', 'cuda/12.1', 'gromacs/2023.3'],
  tpr: 'md.tpr', deffnm: 'md', maxh: 23.5
});

warnings.forEach(w => console.warn(`[${w.level}] ${w.message}`));
\end{lstlisting}

The same four steps performed through the browser produce identical numbers,
because they call the same functions; the difference is only that the browser
supplies the file through \texttt{FileReader} and renders the result. Where a
step produces a figure, the tool also emits a matplotlib script that regenerates
it from the original data file.

% ---------------------------------------------------------------------------
\section{Impact}
\label{sec:impact}

The most direct effect of the architecture is on work that cannot currently be
done at all with server-based tools. Analyses of unpublished results, of
material under a confidentiality agreement, and of data held under an
institutional policy that forbids third-party upload are all outside the reach
of a web service, and are frequently outside the reach of local scripting too
where the workstation is managed and package installation is blocked. Because no
data leaves the browser, these constraints are satisfied structurally rather
than contractually: there is no server to receive the data and consequently no
retention policy to evaluate. For researchers in that position the relevant
comparison is not with a faster tool but with performing the calculation by
hand.

For research questions already being pursued, the improvement is in
reproducibility and in the reliability of the numbers. Extracting the
computation into modules that run unchanged under Node.js means that an analysis
first explored by clicking can be captured as a script, version-pinned and
placed under continuous integration without being rewritten in another
language --- the step at which reproducibility is most often lost. The emitted
matplotlib scripts close the same loop for figures. Independently of that, the
validation strategy raises the floor: every reported statistic is checked
against SciPy or NumPy rather than against the authors' expectations, and the
test suite is intended to be run by users evaluating the software, not solely by
its maintainer, since a library whose numerical claims cannot be checked by its
users offers reproducibility only in principle.

The defects found during extraction (Section~\ref{sec:defects}) carry a wider
lesson, and are the most useful outcome of this work. All three had survived in
software that ran without error and produced plausible numbers. None would have
been caught by testing the code against expectations derived from the same
source; each surfaced only when results were compared with an independent
implementation. A 15\% error in a reported skewness is not visible by
inspection, and neither is a $p$-value floored at zero, which looks like a very
strong result. Authors of scientific software --- particularly of the small,
in-house analysis code that is rarely reviewed --- would be well advised to
check their numerics against an independent implementation before reporting
output.

In daily practice, the change is the removal of a setup step from tasks that
previously carried one. Checking whether an RMSD has converged, or whether a
replicate difference survives an assumption check, becomes an operation
performed at the moment the question arises rather than one deferred until an
environment is available. The tools are also usable in teaching and in shared
computing environments where installation is not an option, since a URL is the
only prerequisite.

STEMKit is released under the MIT licence, hosted at
\url{https://stemkit.net}, developed in the open at
\url{https://github.com/LD-Shell/stemkit} and archived on
Zenodo~\cite{stemkit2026} under the persistent identifier
\url{https://doi.org/10.5281/zenodo.21543112}, which resolves to the current
release. The first archived release is contemporaneous with this manuscript, so
adoption data --- downloads, unique users, citing publications --- does not yet
exist. The arguments above are accordingly claims about what the design makes
possible, not evidence of uptake, and should be read as such. The software is
not used in a commercial setting and has not led to a spin-off company; the
client-side design does, however, make it directly applicable to industrial
research groups whose data-handling policies preclude uploading results to a
third-party service.

% ---------------------------------------------------------------------------
\section{Conclusions}
\label{sec:conclusions}

STEMKit is a suite of 21 browser-based tools for computational chemistry built
on a tested, dependency-injected JavaScript core of 16 modules. The architecture
addresses three constraints simultaneously that existing tooling addresses only
in pairs: it requires no installation, transmits no data and remains scriptable.

The privacy property is structural rather than contractual, because computation
occurs entirely within the browser and there is no server to receive data. The
reproducibility property follows from the separation of computation from
presentation: the identical modules that execute behind the browser interface
execute under Node.js, so an interactive analysis may be captured as a
version-pinned script without rewriting. The 1075-test suite validates every
numerical result against an independent reference, and the three defects that
this strategy exposed --- deflated standardised moments, tail probabilities
floored at zero, and misassigned elements in metalloproteins --- are documented,
corrected and covered by regression tests. Future work will extend format
coverage in the trajectory and structure modules and replace the two log-space
fits with a Levenberg--Marquardt implementation on untransformed data.

% ---------------------------------------------------------------------------
\section*{Acknowledgements}

I thank the maintainers of jStat, Papa Parse, regression.js, Plotly, KaTeX and
3Dmol.js, whose libraries STEMKit builds upon.

\section*{Declaration of competing interest}

The author declares no competing financial or non-financial interests.

%% References. The bibliography is typed inline so the manuscript compiles in a
%% single pdflatex pass. paper.bib holds the same entries for a bibtex build,
%% which uses \bibliographystyle{elsarticle-num} and \bibliography{paper} in
%% place of the thebibliography environment below. Note that paper.bib keys the
%% MDAnalysis entry as michaud-agrawal2011mdanalysis.

\begin{thebibliography}{00}

\bibitem{abraham2015}
M.~J. Abraham, T.~Murtola, R.~Schulz, S.~P\'{a}ll, J.~C. Smith, B.~Hess, and
E.~Lindahl,
\newblock GROMACS: High performance molecular simulations through multi-level
  parallelism from laptops to supercomputers,
\newblock \emph{SoftwareX} \textbf{1--2}, 19--25 (2015).
\newblock \url{https://doi.org/10.1016/j.softx.2015.06.001}

\bibitem{thompson2022}
A.~P. Thompson, H.~M. Aktulga, R.~Berger, D.~S. Bolintineanu, W.~M. Brown,
P.~S. Crozier, P.~J. in~'t~Veld, A.~Kohlmeyer, S.~G. Moore, T.~D. Nguyen,
R.~Shan, M.~J. Stevens, J.~Tranchida, C.~Trott, and S.~J. Plimpton,
\newblock LAMMPS --- a flexible simulation tool for particle-based materials
  modeling at the atomic, meso, and continuum scales,
\newblock \emph{Comput. Phys. Commun.} \textbf{271}, 108171 (2022).
\newblock \url{https://doi.org/10.1016/j.cpc.2021.108171}

\bibitem{tribello2014}
G.~A. Tribello, M.~Bonomi, D.~Branduardi, C.~Camilloni, and G.~Bussi,
\newblock PLUMED 2: New feathers for an old bird,
\newblock \emph{Comput. Phys. Commun.} \textbf{185}(2), 604--613 (2014).
\newblock \url{https://doi.org/10.1016/j.cpc.2013.09.018}

\bibitem{michaudagrawal2011}
N.~Michaud-Agrawal, E.~J. Denning, T.~B. Woolf, and O.~Beckstein,
\newblock MDAnalysis: A toolkit for the analysis of molecular dynamics
  simulations,
\newblock \emph{J. Comput. Chem.} \textbf{32}(10), 2319--2327 (2011).
\newblock \url{https://doi.org/10.1002/jcc.21787}

\bibitem{virtanen2020}
P.~Virtanen, R.~Gommers, T.~E. Oliphant, M.~Haberland, T.~Reddy,
D.~Cournapeau, \emph{et al.},
\newblock SciPy 1.0: fundamental algorithms for scientific computing in Python,
\newblock \emph{Nat. Methods} \textbf{17}(3), 261--272 (2020).
\newblock \url{https://doi.org/10.1038/s41592-019-0686-2}

\bibitem{harris2020}
C.~R. Harris, K.~J. Millman, S.~J. van~der~Walt, R.~Gommers, P.~Virtanen,
D.~Cournapeau, \emph{et al.},
\newblock Array programming with NumPy,
\newblock \emph{Nature} \textbf{585}(7825), 357--362 (2020).
\newblock \url{https://doi.org/10.1038/s41586-020-2649-2}

\bibitem{tiesinga2021}
E.~Tiesinga, P.~J. Mohr, D.~B. Newell, and B.~N. Taylor,
\newblock CODATA recommended values of the fundamental physical constants: 2018,
\newblock \emph{Rev. Mod. Phys.} \textbf{93}(2), 025010 (2021).
\newblock \url{https://doi.org/10.1103/RevModPhys.93.025010}

\bibitem{welch1947}
B.~L. Welch,
\newblock The generalization of `Student's' problem when several different
  population variances are involved,
\newblock \emph{Biometrika} \textbf{34}(1--2), 28--35 (1947).
\newblock \url{https://doi.org/10.1093/biomet/34.1-2.28}

\bibitem{dagostino1971}
R.~B. D'Agostino,
\newblock An omnibus test of normality for moderate and large size samples,
\newblock \emph{Biometrika} \textbf{58}(2), 341--348 (1971).
\newblock \url{https://doi.org/10.1093/biomet/58.2.341}

\bibitem{brown1974}
M.~B. Brown and A.~B. Forsythe,
\newblock Robust tests for the equality of variances,
\newblock \emph{J. Am. Stat. Assoc.} \textbf{69}(346), 364--367 (1974).
\newblock \url{https://doi.org/10.1080/01621459.1974.10482955}

\bibitem{holm1979}
S.~Holm,
\newblock A simple sequentially rejective multiple test procedure,
\newblock \emph{Scand. J. Stat.} \textbf{6}(2), 65--70 (1979).

\bibitem{iglewicz1993}
B.~Iglewicz and D.~C. Hoaglin,
\newblock \emph{How to Detect and Handle Outliers},
\newblock ASQC Quality Press, Milwaukee, WI (1993).

\bibitem{grubbs1969}
F.~E. Grubbs,
\newblock Procedures for detecting outlying observations in samples,
\newblock \emph{Technometrics} \textbf{11}(1), 1--21 (1969).
\newblock \url{https://doi.org/10.1080/00401706.1969.10490657}

\bibitem{press2007}
W.~H. Press, S.~A. Teukolsky, W.~T. Vetterling, and B.~P. Flannery,
\newblock \emph{Numerical Recipes: The Art of Scientific Computing},
\newblock 3rd ed., Cambridge University Press, Cambridge (2007).

\bibitem{stemkit2026}
O.~M. Daramola,
\newblock STEMKit: A client-side toolkit for computational chemistry and
  scientific data analysis (v0.1.0), Zenodo (2026).
\newblock \url{https://doi.org/10.5281/zenodo.21543112}

\end{thebibliography}

%% The SoftwareX template Version 6 (March 2026) no longer carries the
%% "Current executable software version" (S1-S8) table that earlier versions
%% included, so it is omitted here.

\end{document}
\endinput
